% !TEX root = ../main.tex
\begin{ouverture}[photons]{20}{Lumière : ondes et photons}{Partie 4 · Ondes et signaux}
{Ondes électromagnétiques\enspace·\enspace Domaines spectraux\enspace·\enspace Photon\enspace·\enspace Niveaux d'énergie\enspace·\enspace Spectres de raies}
\objectif{Utiliser une échelle de fréquences ou de longueurs d'onde pour identifier un domaine spectral ; citer des ordres de grandeur (imagerie médicale, visible, wifi, micro-ondes).}
\objectif{Exploiter la relation $\lambda=\dfrac{c}{\nu}$ entre longueur d'onde, célérité de la lumière et fréquence.}
\objectif{Utiliser l'expression de l'énergie d'un photon, $E=h\,\nu$, en joules et en électronvolts.}
\objectif{Exploiter un diagramme de niveaux d'énergie pour interpréter un spectre d'émission ou d'absorption.}
\prerequis{Spectres d'émission et d'absorption, spectres de raies, lumière blanche (classe de seconde).}
\prerequis{Période, fréquence et longueur d'onde (chapitre 17) ; couleurs (chapitre 19) ; puissances de 10.}
\end{ouverture}
\essentiel
\essSection{1}{Les ondes électromagnétiques}
\begin{definition}[Onde électromagnétique]
La lumière est une \textbf{onde électromagnétique} : contrairement aux ondes mécaniques, elle se propage aussi dans le vide, à la célérité
$c=\qty{3,00e8}{m.s^{-1}}$. Une radiation est caractérisée par sa \textbf{fréquence} $\nu$ (en Hz), fixée par la source, ou par sa \textbf{longueur
d'onde dans le vide} $\lambda$ (en m) :
\[\lambda=\frac{c}{\nu}=c\,T.\]
Une radiation de \qty{500}{nm} a ainsi une fréquence $\nu=\dfrac{c}{\lambda}=\qty{6,00e14}{Hz}$.
\end{definition}
\par\penalty0
\schema{20-domaines}
\begin{propriete}[Domaines et ordres de grandeur]
Les ondes électromagnétiques sont classées en domaines, aux limites approximatives : rayons $\gamma$, rayons X (imagerie médicale, de
\qty{0,01}{nm} à \qty{10}{nm}), ultraviolet (de \qty{10}{nm} à \qty{400}{nm}), \textbf{visible} (de \qty{400}{nm} à \qty{800}{nm}), infrarouge (de \qty{800}{nm} à \qty{1}{mm}),
micro-ondes (four, wifi à \qty{2,4}{GHz} : $\lambda=\qty{12,5}{cm}$) et ondes radio (radio FM à \qty{100}{MHz} : $\lambda=\qty{3,0}{m}$).
\end{propriete}
\begin{savoirfaire}[Identifier le domaine d'une radiation]
\begin{enumerate}
\item Exprimer la grandeur donnée en unités SI : $\lambda$ en mètres (\qty{1}{nm}~$=\qty{e-9}{m}$), $\nu$ en hertz (\qty{1}{GHz}~$=\qty{e9}{Hz}$).
\item Si la fréquence est donnée, calculer la longueur d'onde $\lambda=\dfrac{c}{\nu}$.
\item Situer $\lambda$ sur l'échelle des domaines, en comparant d'abord aux limites du visible, \qty{400}{nm} et \qty{800}{nm}.
\end{enumerate}
Exemple : un radar émet à \qty{10}{GHz} ; $\lambda=\dfrac{\num{3,00e8}}{\num{1,0e10}}=\qty{3,0e-2}{m}$, soit \qty{3,0}{cm} : ce sont des micro-ondes.
\end{savoirfaire}
\essSection{2}{Le photon}
\begin{definition}[Modèle particulaire de la lumière]
Pour décrire les échanges d'énergie avec la matière, on modélise la lumière par un flux de particules, les \textbf{photons}. Un photon a
une masse nulle, se déplace à la célérité $c$ et transporte l'énergie
\[E=h\,\nu=\frac{h\,c}{\lambda},\]
où $h=\qty{6,63e-34}{J.s}$ est la constante de Planck. À l'échelle de l'atome, on utilise l'\textbf{électronvolt} : $\qty{1}{eV}=\qty{1,60e-19}{J}$.
\end{definition}
\begin{propriete}[Ordres de grandeur]
Un photon de \qty{500}{nm} transporte
\[E=\frac{h\,c}{\lambda}=\frac{\num{6,63e-34}\times\num{3,00e8}}{\num{500e-9}}=\qty{3,98e-19}{J},\quad\text{soit}\quad\frac{\num{3,98e-19}}{\num{1,60e-19}}=\qty{2,49}{eV}.\]
Les photons visibles transportent de \qty{1,6}{eV} (rouge) à \qty{3,1}{eV} (violet) ; ceux des rayons X, de l'ordre de \qty{e4}{eV} ; ceux des
micro-ondes, de l'ordre de \qty{e-5}{eV}.
\end{propriete}
\begin{attention}
L'énergie d'un photon varie en sens inverse de sa longueur d'onde : un photon ultraviolet est plus énergétique qu'un photon rouge,
c'est pourquoi les UV abîment la peau. Dans $E=\dfrac{h\,c}{\lambda}$, $\lambda$ s'exprime en mètres et $E$ en joules : convertir d'abord les nm et les eV.
\end{attention}
\essSection{3}{Niveaux d'énergie et échanges lumière-matière}
\begin{loi}[Quantification de l'énergie des atomes]
L'énergie d'un atome est \textbf{quantifiée} : elle ne peut prendre que certaines valeurs, les \textbf{niveaux d'énergie}. Le plus bas est
l'\textbf{état fondamental}, les autres sont des \textbf{états excités}. Lors d'une transition entre deux niveaux, l'atome émet ou absorbe un
photon dont l'énergie est égale à l'écart entre ces niveaux :
\[\Delta E=E_\text{haut}-E_\text{bas}=h\,\nu=\frac{h\,c}{\lambda}.\]
\textbf{Émission} : l'atome descend d'un niveau et émet le photon. \textbf{Absorption} : l'atome monte d'un niveau en absorbant un photon,
seulement si l'énergie de ce photon est exactement égale à l'écart entre deux niveaux.
\end{loi}
\par\penalty0
\schema{20-niveaux}
\begin{python}
h, c, e = 6.63e-34, 3.00e8, 1.60e-19
for n in [3, 4, 5, 6]:
    dE = 13.6 / 2**2 - 13.6 / n**2       # eV : ecart entre les niveaux n et 2
    lam = h * c / (dE * e)                # m
    print(n, '-> 2 :', round(dE, 2), 'eV ;', round(lam * 1e9), 'nm')
\end{python}
\begin{console}
3 -> 2 : 1.89 eV ; 658 nm
4 -> 2 : 2.55 eV ; 488 nm
5 -> 2 : 2.86 eV ; 435 nm
6 -> 2 : 3.02 eV ; 411 nm
\end{console}
\noindent\capacite{Capacité numérique} Les niveaux de l'atome d'hydrogène valent $E_n=-\dfrac{\qty{13,6}{eV}}{n^2}$. Le programme calcule les raies émises quand l'atome
redescend sur le niveau $n=2$ : ce sont les quatre raies visibles de l'hydrogène, mesurées à 656, 486, 434 et \qty{410}{nm} ; le faible écart
(\qty{0,3}{\percent}) vient de l'arrondi des constantes.
\essSection{4}{Spectres de raies}
\begin{propriete}[Spectres d'émission et d'absorption]
Un gaz chaud émet un \textbf{spectre de raies} : chaque raie correspond à une transition vers un niveau inférieur. Une lumière blanche qui
traverse un gaz perd les radiations qui font monter ses atomes d'un niveau à un autre : le spectre d'absorption présente des raies noires
\textbf{aux mêmes longueurs d'onde} que les raies d'émission. Propres à chaque élément, ces raies en sont la \textbf{signature} : elles révèlent
la composition des étoiles.
\end{propriete}
\schema{20-spectres}
\begin{experience}[Le spectre d'une lampe spectrale]
On observe avec un spectroscope, ou l'on enregistre avec un spectromètre relié à un ordinateur, la lumière d'une lampe à vapeur de
mercure ou de sodium. On mesure la longueur d'onde des raies, on calcule l'énergie des photons correspondants et l'on retrouve les
transitions sur le diagramme de niveaux fourni.
\end{experience}
\begin{rediger}
Lire les niveaux sur le diagramme, calculer $\Delta E$ en eV puis le convertir en joules, appliquer $\lambda=\dfrac{h\,c}{\Delta E}$, donner $\lambda$ en nm avec trois
chiffres significatifs et situer la radiation dans son domaine (UV, visible, infrarouge).
\end{rediger}
\exercices

\rubrique{Découverte}
\begin{blocLivre}{ex-20-1}
\begin{exercice}[1]{Des ondes au quotidien}
On considère quatre rayonnements : (a) le signal d'une box wifi, de fréquence \qty{2,4}{GHz} ; (b) les micro-ondes d'un four, de fréquence
\qty{2,45}{GHz} ; (c) les rayons d'un appareil de radiographie, de longueur d'onde \qty{0,10}{nm} ; (d) le rayonnement capté par une caméra
thermique, de longueur d'onde \qty{10}{\micro m}.
\begin{questions}
\item Calculer la longueur d'onde dans le vide des rayonnements (a) et (b).
\item Calculer la fréquence des rayonnements (c) et (d).
\item À l'aide de l'échelle des domaines de L'essentiel, nommer le domaine de chaque rayonnement.
\item Lequel a la plus grande fréquence ? La plus grande longueur d'onde ?
\end{questions}
\end{exercice}
\end{blocLivre}
\begin{blocLivre}{ex-20-2}
\begin{exercice}[1]{Photons rouge et violet}
\begin{donnees}
$h=\qty{6,63e-34}{J.s}$ ; $c=\qty{3,00e8}{m.s^{-1}}$ ; $\qty{1}{eV}=\qty{1,60e-19}{J}$.
\end{donnees}
\begin{questions}
\item Calculer l'énergie d'un photon de \qty{700}{nm} (rouge), en joules puis en électronvolts.
\item Même question pour un photon de \qty{400}{nm} (violet).
\item Lequel est le plus énergétique ? Combien de fois plus ?
\item Classer par énergie croissante un photon X, un photon infrarouge, un photon ultraviolet et un photon visible.
\end{questions}
\end{exercice}
\end{blocLivre}
\begin{blocLivre}{ex-20-3}
\begin{exercice}[1]{Le diagramme de l'hydrogène}
\schema{20-ex03}
\begin{donnees}
$h=\qty{6,63e-34}{J.s}$ ; $c=\qty{3,00e8}{m.s^{-1}}$ ; $\qty{1}{eV}=\qty{1,60e-19}{J}$.
\end{donnees}
\begin{questions}
\item Quel niveau correspond à l'état fondamental ? Que représente le niveau d'énergie 0 ?
\item L'atome passe du niveau $n=2$ au niveau $n=1$. Émet-il ou absorbe-t-il un photon ? Calculer l'énergie de ce photon, en eV puis en J.
\item Calculer la longueur d'onde correspondante. Dans quel domaine se situe-t-elle ?
\item Même question pour la transition du niveau $n=3$ au niveau $n=1$.
\end{questions}
\end{exercice}
\end{blocLivre}
\begin{blocLivre}{ex-20-4}
\begin{exercice}[1]{La lampe au sodium}
Les anciens lampadaires au sodium émettent une lumière jaune-orangé presque monochromatique, de longueur d'onde \qty{589}{nm}.
\begin{donnees}
$h=\qty{6,63e-34}{J.s}$ ; $c=\qty{3,00e8}{m.s^{-1}}$ ; $\qty{1}{eV}=\qty{1,60e-19}{J}$ ; deux niveaux de l'atome de sodium : $E_0=\qty{-5,14}{eV}$ et
$E_1=\qty{-3,03}{eV}$.
\end{donnees}
\begin{questions}
\item Calculer la fréquence de cette radiation.
\item Calculer l'énergie d'un photon émis, en joules puis en électronvolts.
\item Montrer que cette émission correspond à une transition entre les niveaux $E_1$ et $E_0$.
\item Sous ces lampadaires, une voiture bleue paraît noire. Expliquer.
\end{questions}
\end{exercice}
\end{blocLivre}
\begin{blocLivre}{ex-20-5}
\begin{exercice}[1]{Micro-ondes et ultraviolets}
Un four émet des micro-ondes de fréquence \qty{2,45}{GHz} ; le Soleil émet notamment des ultraviolets de \qty{300}{nm}.
\begin{donnees}
$h=\qty{6,63e-34}{J.s}$ ; $c=\qty{3,00e8}{m.s^{-1}}$ ; $\qty{1}{eV}=\qty{1,60e-19}{J}$.
\end{donnees}
\begin{questions}
\item Calculer l'énergie d'un photon du four, en joules puis en électronvolts.
\item Calculer l'énergie d'un photon ultraviolet de \qty{300}{nm}, en électronvolts.
\item Combien de fois le photon ultraviolet est-il plus énergétique ?
\item Il faut quelques électronvolts pour casser une liaison dans une molécule d'ADN. Expliquer pourquoi les UV sont dangereux pour la peau
alors que les micro-ondes ne peuvent pas casser ces liaisons.
\end{questions}
\end{exercice}
\end{blocLivre}
\rubrique{Entraînement}
\begin{blocLivre}{ex-20-6}
\begin{exercice}[2]{Les raies de l'hydrogène, en Python}
\begin{python}
h, c, e = 6.63e-34, 3.00e8, 1.60e-19
def domaine(lam):
    if lam < 400e-9:
        return 'UV'
    if lam <= 800e-9:
        return 'visible'
    return 'IR'
for haut, bas in [(2, 1), (3, 2), (4, 3)]:
    dE = 13.6 / bas**2 - 13.6 / haut**2       # eV
    lam = h * c / (dE * e)                    # m
    print(haut, '->', bas, ':', round(lam * 1e9), 'nm', domaine(lam))
\end{python}
\begin{console}
2 -> 1 : 122 nm UV
3 -> 2 : 658 nm visible
4 -> 3 : 1880 nm IR
\end{console}
\begin{questions}
\item \capacite{Capacité numérique}Les niveaux de l'hydrogène valent $E_n=-\dfrac{\qty{13,6}{eV}}{n^2}$. Montrer que \code{dE} est l'énergie du photon émis lors de la transition.
\item Retrouver par le calcul la longueur d'onde affichée pour la transition $2\to1$.
\item L'atome d'hydrogène n'émet-il que de la lumière visible ? Justifier.
\item Que faut-il ajouter au programme pour étudier la transition $5\to2$ ? Prévoir la ligne affichée.
\end{questions}
\end{exercice}
\end{blocLivre}
\begin{blocLivre}{ex-20-7}
\begin{exercice}[2]{Le pointeur laser}
Un pointeur laser rouge émet une radiation de \qty{650}{nm} avec une puissance de \qty{1,0}{mW}, c'est-à-dire \qty{1,0e-3}{J} d'énergie lumineuse par seconde.
\begin{donnees}
$h=\qty{6,63e-34}{J.s}$ ; $c=\qty{3,00e8}{m.s^{-1}}$ ; $\qty{1}{eV}=\qty{1,60e-19}{J}$.
\end{donnees}
\begin{questions}
\item Calculer la fréquence de la radiation et l'énergie d'un photon, en joules et en électronvolts.
\item Combien de photons le laser émet-il chaque seconde ?
\item Un laser vert de \qty{532}{nm} a la même puissance. Émet-il plus ou moins de photons par seconde ? Justifier par un calcul.
\item Quelle énergie le laser rouge émet-il en une minute ? Pourquoi est-il pourtant dangereux de le diriger vers un œil ?
\end{questions}
\end{exercice}
\end{blocLivre}
\begin{blocLivre}{ex-20-8}
\begin{exercice}[2]{La lampe à vapeur de mercure}
\schema{20-ex08}
\begin{donnees}
$h=\qty{6,63e-34}{J.s}$ ; $c=\qty{3,00e8}{m.s^{-1}}$ ; $\qty{1}{eV}=\qty{1,60e-19}{J}$.
\end{donnees}
\begin{questions}
\item Calculer la longueur d'onde des raies émises lors des transitions $E_4\to E_1$ et $E_4\to E_2$. Quelle est leur couleur ?
\item Calculer la longueur d'onde de la raie émise lors de la transition $E_1\to E_0$. Dans quel domaine se situe-t-elle ?
\item Dans un tube fluorescent, la vapeur de mercure émet surtout cette dernière raie ; la paroi est recouverte d'une poudre qui la
transforme en lumière visible. Pourquoi cette poudre est-elle indispensable ?
\item Un atome de mercure dans son état fondamental reçoit un photon de \qty{4,00}{eV}, puis un photon de \qty{4,89}{eV}. Lequel peut-il absorber ?
\end{questions}
\end{exercice}
\end{blocLivre}
\begin{blocLivre}{ex-20-9}
\begin{exercice}[2]{La radiographie}
Un appareil de radiographie émet des rayons X de longueur d'onde \qty{0,050}{nm}.
\begin{donnees}
$h=\qty{6,63e-34}{J.s}$ ; $c=\qty{3,00e8}{m.s^{-1}}$ ; $\qty{1}{eV}=\qty{1,60e-19}{J}$.
\end{donnees}
\begin{questions}
\item Calculer la fréquence de ces rayons X.
\item Calculer l'énergie d'un photon X, en joules puis en électronvolts.
\item Combien de fois ce photon est-il plus énergétique qu'un photon visible de \qty{500}{nm} ?
\item Les rayons X traversent les tissus mous mais sont arrêtés par les os. Comment obtient-on l'image ? Pourquoi le manipulateur se
protège-t-il à chaque cliché ?
\end{questions}
\end{exercice}
\end{blocLivre}
\begin{blocLivre}{ex-20-10}
\begin{exercice}[2]{Absorption par l'hydrogène}
On utilise le diagramme de niveaux de l'hydrogène de l'exercice 3 ; l'atome est dans son état fondamental.
\begin{donnees}
$h=\qty{6,63e-34}{J.s}$ ; $c=\qty{3,00e8}{m.s^{-1}}$ ; $\qty{1}{eV}=\qty{1,60e-19}{J}$.
\end{donnees}
\begin{questions}
\item On envoie des photons de \qty{10,2}{eV}, \qty{11,0}{eV} et \qty{12,1}{eV}. Lesquels peuvent être absorbés ? Justifier.
\item Quelle énergie minimale faut-il fournir pour ioniser l'atome ? Calculer la longueur d'onde du photon correspondant.
\item Un photon de \qty{15,0}{eV} est absorbé. Que devient l'atome ? Avec quelle énergie cinétique l'électron s'éloigne-t-il ?
\item De l'hydrogène à température ambiante, éclairé en lumière blanche, n'absorbe aucune radiation visible. Expliquer.
\end{questions}
\end{exercice}
\end{blocLivre}
\begin{blocLivre}{ex-20-11}
\begin{exercice}[2]{Les DEL de couleur}
Dans une diode électroluminescente (DEL), chaque électron qui traverse le composant peut émettre un photon ; l'énergie de ce photon est
fixée par le matériau semi-conducteur.
\begin{donnees}
$h=\qty{6,63e-34}{J.s}$ ; $c=\qty{3,00e8}{m.s^{-1}}$ ; $\qty{1}{eV}=\qty{1,60e-19}{J}$.
\end{donnees}
\begin{questions}
\item Une DEL émet des photons de \qty{1,9}{eV}. Calculer la longueur d'onde émise et en déduire sa couleur.
\item Même question pour une DEL dont les photons ont une énergie de \qty{2,76}{eV}.
\item Une DEL verte émet à \qty{525}{nm}. Quelle est l'énergie de ses photons ?
\item Une DEL « blanche » est une DEL bleue recouverte d'une poudre qui transforme une partie du bleu en jaune. Pourquoi la lumière
émise paraît-elle blanche ?
\end{questions}
\end{exercice}
\end{blocLivre}
\rubrique{Objectif DS}
\begin{blocLivre}{ex-20-12}
\begin{exercice}[3]{Le four à micro-ondes}
\begin{docu}{1}{Le magnétron}
Le magnétron d'un four émet des micro-ondes de fréquence \qty{2,45}{GHz} avec une puissance de \qty{800}{W}. Elles sont absorbées par les
molécules d'eau des aliments, qui s'échauffent.
\end{docu}
\begin{docu}{2}{Une expérience}
On retire le plateau tournant et l'on chauffe une tablette de chocolat. Les zones fondues sont séparées d'une demi-longueur d'onde.
\schema{20-ex12}
\end{docu}
\begin{donnees}
$h=\qty{6,63e-34}{J.s}$ ; $c=\qty{3,00e8}{m.s^{-1}}$.
\end{donnees}
\begin{questions}
\item Calculer la longueur d'onde des micro-ondes. Vérifier qu'elles appartiennent bien au domaine des micro-ondes.
\item Mesurer, le plus précisément possible, la distance entre deux zones fondues voisines. En déduire une valeur de la célérité $c$.
\item \initiative Combien de photons le magnétron émet-il chaque seconde ?
\item Expliquer le rôle du plateau tournant.
\end{questions}
\end{exercice}
\end{blocLivre}
\begin{blocLivre}{ex-20-13}
\begin{exercice}[3]{Les feux d'artifice}
\begin{docu}{1}{Des sels qui colorent les flammes}
Les couleurs des feux d'artifice proviennent de sels métalliques chauffés : les sels de sodium émettent surtout à \qty{589}{nm}, ceux de
strontium vers \qty{650}{nm} et ceux de baryum vers \qty{520}{nm}.
\end{docu}
\begin{docu}{2}{Excitation et désexcitation}
Dans la flamme, les chocs avec les autres particules portent les atomes dans des états excités ; ils redescendent ensuite vers des
niveaux plus bas en émettant de la lumière.
\end{docu}
\begin{donnees}
$h=\qty{6,63e-34}{J.s}$ ; $c=\qty{3,00e8}{m.s^{-1}}$ ; $\qty{1}{eV}=\qty{1,60e-19}{J}$.
\end{donnees}
\begin{questions}
\item Le spectre de la lumière émise est-il un spectre d'émission ou d'absorption ? Continu ou de raies ? Justifier.
\item Indiquer la couleur produite par chaque sel et calculer l'énergie des photons correspondants, en électronvolts.
\item \initiative Un artificier veut une gerbe jaune sans utiliser de sodium. Quels sels peut-il mélanger ? Justifier.
\item Un spectateur porte un pull bleu. De quelle couleur paraît-il lorsque la scène n'est éclairée que par une gerbe au sodium ?
\end{questions}
\end{exercice}
\end{blocLivre}
\begin{blocLivre}{ex-20-14}
\begin{exercice}[3]{La cellule photovoltaïque}
\begin{docu}{1}{Le silicium}
Dans le silicium d'une cellule photovoltaïque, un photon peut libérer un électron et produire du courant si son énergie est au moins
égale à \qty{1,12}{eV}. Chaque électron libéré ne fournit que \qty{1,12}{eV} ; l'excédent d'énergie du photon est perdu sous forme de chaleur.
\end{docu}
\begin{docu}{2}{La lumière du Soleil}
Le rayonnement solaire reçu au sol s'étend de \qty{300}{nm} à \qty{2500}{nm} environ ; près de la moitié de son énergie est transportée par
l'infrarouge.
\end{docu}
\begin{donnees}
$h=\qty{6,63e-34}{J.s}$ ; $c=\qty{3,00e8}{m.s^{-1}}$ ; $\qty{1}{eV}=\qty{1,60e-19}{J}$.
\end{donnees}
\begin{questions}
\item Calculer la longueur d'onde maximale des photons utilisables. Dans quel domaine se situe-t-elle ?
\item Un photon de \qty{500}{nm} est absorbé. Calculer l'énergie perdue en chaleur et la fraction de son énergie qu'elle représente.
\item Un photon infrarouge de \qty{1500}{nm} peut-il produire du courant ?
\item \initiative Expliquer pourquoi le rendement d'une cellule au silicium reste bien inférieur à \qty{100}{\percent}.
\end{questions}
\end{exercice}
\end{blocLivre}
\begin{blocLivre}{ex-20-15}
\begin{exercice}[3]{Les raies de Fraunhofer}
\begin{docu}{1}{Le spectre du Soleil}
En 1814, Joseph von Fraunhofer observe des centaines de raies sombres dans le spectre de la lumière solaire.
\schema{20-ex15}
\end{docu}
\begin{docu}{2}{L'hydrogène}
Les niveaux d'énergie de l'atome d'hydrogène valent $E_n=-\dfrac{\qty{13,6}{eV}}{n^2}$, avec $n$ entier.
\end{docu}
\begin{donnees}
$h=\qty{6,63e-34}{J.s}$ ; $c=\qty{3,00e8}{m.s^{-1}}$ ; $\qty{1}{eV}=\qty{1,60e-19}{J}$.
\end{donnees}
\begin{questions}
\item Ce spectre est-il un spectre d'émission ou d'absorption ? Quelle est l'origine des raies sombres ?
\item Calculer l'énergie des photons de longueurs d'onde \qty{656,3}{nm} et \qty{486,1}{nm}.
\item \initiative Montrer que ces deux raies révèlent la présence d'hydrogène dans l'atmosphère du Soleil.
\item La raie à \qty{589,0}{nm} est due à un élément rencontré dans ce chapitre. Lequel ?
\end{questions}
\end{exercice}
\end{blocLivre}
